% ===================== Chapter 10: File Input/Output (ARC500) =====================
% Requires your project preamble to define:
%   \h{...}, \hh{...}, \hhh{...}
%   \code{<label>}{python}{<code>}{<caption>}
% All Python examples below use EXACTLY 4 spaces per indent (no tabs).

\h{File Input/Output}

In programming, variables hold data temporarily; when your script finishes, that data vanishes. \emph{File Input/Output} (or \emph{File I/O}) is the bridge between your program and persistent data stored on your computer's disk. \c{Input} is the act of reading data \c{from} a file into your script's memory. \c{Output} is the act of writing data \c{to} a file, saving it for later use by you, your colleagues, or other software.

This single capability "saving and loading data" is what elevates a simple script into a powerful and reusable tool. It makes your work \c{reproducible, shareable, and automated}.

\hh{Why File I/O is Essential}

Think about the data that flows through a typical architecture project: schedules, material specifications, building codes, cost estimates, and analysis results. Managing this information manually is slow and prone to error. With Python, you can automate these workflows.

For example, you could write a script to:
\begin{itemize}
  \item \textbf{Import data for analysis:} Read a CSV file of a room area program to automatically check if the design meets the required totals.
  \item \textbf{Generate reports:} Export a neatly formatted text file or spreadsheet detailing material quantities and costs, calculated directly from a model's data.
  \item \textbf{Create logs:} Keep a running record of a generative design process, logging the key parameters and results of each iteration to create a traceable audit trail.
  \item \textbf{Link software tools:} Pass structural analysis results from one application into a spreadsheet for visualization, or take quantities from a BIM model and format them for a specification document.
  \item \textbf{Manage project files:} Write a script to organize thousands of project photos or automatically back up critical drawing files into a dated archive folder.
\end{itemize}

Ultimately, good file I/O practices enable you to:
\begin{itemize}
  \item \textbf{Automate} tedious data entry and formatting tasks.
  \item \textbf{Validate} your work with scripts that can check data for errors or inconsistencies.
  \item \textbf{Standardize} how data is exchanged between team members and different software, avoiding clumsy copy-pasting.
  \item \textbf{Scale} your work to handle large datasets, such as processing sensor readings from a building or analyzing extensive climate data for a site.
\end{itemize}

\begin{tcolorbox}[
    title=Focus of this chapter,
    colframe=orange,
    colback=white,
    arc=2mm
  ]
  \begin{itemize}
    \item Open and close files safely using Python's context manager (\c{with open(...)}).
    \item Understand file \textbf{modes} like read (\c{"r"}), write (\c{"w"}), and append (\c{"a"}).
    \item Read and write text content line-by-line or all at once.
    \item Navigate your computer's filesystem using the modern \c{pathlib} library to easily find files and create folders.
    \item Handle common, structured data formats like CSV and JSON using Python's built-in libraries.
    \item Apply these skills to practical tasks, like generating reports and creating append-only logs for your design experiments.
  \end{itemize}
\end{tcolorbox}
% ====================================================================
% PART I — FILE I/O BASICS
% ====================================================================

\hh{Opening and Closing Files: The 'with' Statement}

Think of a file on your disk as a notebook on a shelf. To use it, you first need to get it. In Python, the \c{open()} function retrieves the notebook for you, returning a special \c{file object}.

The safest and most common way to open a file is with a \c{with} block. This structure is a \c{context manager}: it automatically closes the file for you when you are done (i.e., when the indented block ends), even if your code runs into an error. This prevents data corruption and other issues.

\code{c13_tp_open_close}{python}{
    # The 'with... as f:' syntax opens a file and assigns it a temporary name.
    # Here, 'f' is a variable—a "handle" to the file. You could also name it
    # 'outfile', 'report_file', etc.

    # 1. Open a file in "write" mode ('w') to create it.
    with open("hello_arc500.txt", "w", encoding="utf-8") as f:
        f.write("Hello from ARC500!")  # Write a string to the file.

    # The 'with' block has ended, so "hello_arc500.txt" is now closed.

    # 2. Open the same file in "read" mode ('r') to see its contents.
    with open("hello_arc500.txt", "r", encoding="utf-8") as infile:
        content = infile.read()        # Read the entire file into a string.

    print(content)
    # Outcome:
    # Hello from ARC500!
}{Use a `with` block to ensure files are closed automatically.}

% --------------------------------------------------------------------
\hh{File Modes: Stating Your Intent}

When you open a file, you must tell Python what you intend to do. This is called the file \c{mode}. Think of it as deciding whether you will only read a document, write a new one from scratch, or add notes to the end of an existing one.

The most common modes are:
\begin{itemize}
  \item \c{"r"} -- \textbf{Read} (the default). Opens a file for reading. Throws an error if the file doesn't exist.
  \item \c{"w"} -- \textbf{Write}. Opens a file for writing. Creates a new file or completely overwrites an existing one.
  \item \c{"a"} -- \textbf{Append}. Opens a file for writing. Creates a new file if one doesn't exist, but if it does, adds new content to the very end.
  \item \c{"x"} -- \textbf{Exclusive create}. Creates a new file for writing, but fails with an error if the file already exists.
\end{itemize}
You can also add \c{"b"} for binary files (e.g., images, PDFs) or \c{"t"} for text files (the default). For this chapter, we will focus on text files.

\code{c13_tp_modes}{python}{
    # Demonstrate the critical difference between "w" (overwrite) and "a" (append).

    # Use 'w' to create a new report, erasing any previous version.
    with open("report.txt", "w", encoding="utf-8") as f:
        f.write("Project Area Report\n")
        f.write("-------------------\n")
        f.write("Level 1: 1500 m2\n")

    # Later, use 'a' to add more data to the same report without deleting it.
    with open("report.txt", "a", encoding="utf-8") as f:
        f.write("Level 2: 1350 m2\n")

    # Read the final result.
    with open("report.txt", "r", encoding="utf-8") as f:
        print(f.read())
    # Outcome:
    # Project Area Report
    # -------------------
    # Level 1: 1500 m2
    # Level 2: 1350 m2
}{Use "w" to start fresh and "a" to add to an existing file.}

\hhh{A Quick Note: What is {encoding="utf-8"}?}
Think of \c{encoding="utf-8"} as a \textbf{universal legend for text}, similar to the legend on an architectural drawing.

A computer doesn't store the letter 'A' or the symbol '°'; it stores numbers (called bytes). An \textbf{encoding} is the rulebook that translates human-readable characters into these numbers and back again. By specifying \texttt{"utf-8"}, you are choosing the modern, universal standard that can handle nearly every character and symbol from any language in the world, which is essential for technical documents and international collaboration.

\hhh*{Why is this critical?}
If you don't specify an encoding, your computer will guess. If you share a script with a colleague whose computer guesses differently (e.g., between Windows and macOS), any special characters (like \texttt{€}, \texttt{ü}, or \texttt{°}) can become garbled. Using \texttt{utf-8} ensures your files work correctly for everyone, everywhere.

\textbf{Golden rule:} Always specify \c{encoding="utf-8"} when you \c{open()} a text file. It’s a simple habit that makes your code \textbf{reliable} and \textbf{portable}.


% --------------------------------------------------------------------
\hh{Writing Text}

The \c{f.write()} method writes a string to a file. Crucially, it does not automatically add a line break. You must explicitly add the newline character, \c{\n}, to your string wherever you want to start a new line. This gives you complete control over the output format.

\code{c13_tp_write}{python}{
    # A list of tuples containing room data: (Room ID, Area in sq. meters).
    room_data = [("A-101", 18.5), ("A-102", 21.3), ("A-103", 15.2)]

    with open("areas.csv", "w", encoding="utf-8") as f:
        # Write a header row, followed by a newline character.
        f.write("room_id,area_m2\n")
        
        # Loop through the data and write each record as a new line.
        for room_id, area in room_data:
            # Use an f-string to format the data into a CSV row.
            f.write(f"{room_id},{area}\n")

    # Verify the contents of the newly created CSV file.
    with open("areas.csv", "r", encoding="utf-8") as f:
        print(f.read())
    # Outcome:
    # room_id,area_m2
    # A-101,18.5
    # A-102,21.3
    # A-103,15.2
}{The {f.write()} method requires you to add newline characters (\\n) manually.}

% --------------------------------------------------------------------
\hh{Reading Text}

Python offers several ways to read from a text file. The best method depends on the file size and what you need to accomplish.

\begin{enumerate}
  \item \c{f.read()} -- Best for small files. Reads the entire file content into a single string.
  \item \c{f.readlines()} -- Reads all lines into a list of strings. Can use a lot of memory for large files.
  \item \c{f.readline()} -- Reads just one line at a time. Good for processing files line-by-line manually.
  \item \c{for line in f:} -- Best for most cases, especially large files. Reads the file line by line without loading it all into memory at once. This is the most efficient and Pythonic method.
\end{enumerate}

\code{c13_tp_read_readline}{python}{
    # First, create a sample file to work with.
    with open("rooms_demo.txt", "w", encoding="utf-8") as f:
        f.write("B-201,19.0\nB-202,22.4\nB-203,17.6\n")

    # Method 4: Iterate line-by-line to calculate total area (most common pattern).
    total_area = 0.0
    with open("rooms_demo.txt", "r", encoding="utf-8") as f:
        for line in f:
            # Each 'line' from the file includes a trailing '\n', so we strip it.
            clean_line = line.strip()
            if clean_line: # Make sure the line isn't empty
                room_id, area_str = clean_line.split(",")
                total_area += float(area_str)

    print(f"Total Area: {total_area:.2f} m2")
    # Outcome:
    # Total Area: 59.00 m2
}{Iterating with `for line in f:` is the most memory-efficient way to read a file.}

% --------------------------------------------------------------------
\hh{Managing Project Files with pathlib}

Beyond reading and writing, you often need to perform basic file management, like creating folder structures, cleaning up temporary files, or archiving results. The \texttt{pathlib} library is the modern, object-oriented way to handle these tasks. Instead of treating file paths as simple text strings, it creates "Path" objects that have useful properties and methods, making your code cleaner, safer, and compatible across Windows, macOS, and Linux.

\hhh{Creating and Joining Paths}
The foundation of \texttt{pathlib} is the \texttt{Path} object. The most intuitive feature is using the forward slash operator \texttt{/} to build paths, which automatically handles the correct separator for the operating system.

\code{c13_tp_path_creation}{python}{
    from pathlib import Path

    # The / operator is the recommended way to join path components.
    project_dir = Path("P-2025-City-Library")
    drawings_path = project_dir / "02_Drawings" / "CAD"
    output_path = drawings_path / "A-101_Floor-Plan.dwg"
    
    # pathlib makes paths easy to read and manage.
    print(f"Full path: {output_path}")

    # You can also get special directory paths programmatically.
    print(f"Current Directory: {Path.cwd()}")
    print(f"User's Home: {Path.home()}")
    # Outcome (example):
    # Full path: P-2025-City-Library/02_Drawings/CAD/A-101_Floor-Plan.dwg
    # Current Directory: /path/to/your/current/folder
    # User's Home: /Users/your_username
}{Constructing paths with the / operator and accessing special directories.}

\hhh{Extracting Path Components}
A Path object lets you easily access individual parts of the path, which is incredibly useful for sorting, filtering, or renaming files.

\code{c13_tp_path_components}{python}{
    from pathlib import Path

    file_path = Path("project/renders/Exterior_View_03.png")

    print(f"  Parent directory: {file_path.parent}")
    print(f"  Full name: {file_path.name}")
    print(f"  Name without extension (stem): {file_path.stem}")
    print(f"  Extension (suffix): {file_path.suffix}")
    # Outcome:
    #   Parent directory: project/renders
    #   Full name: Exterior_View_03.png
    #   Name without extension (stem): Exterior_View_03
    #   Extension (suffix): .png
}{Accessing the parent, name, stem, and suffix of a path.}

\hhh{Checking Path Properties}
Before you try to read or write a file, you can check its status, for example, to see if it exists or if it's a file versus a directory.

\code{c13_tp_path_checks}{python}{
    from pathlib import Path

    # Let's create a directory and a file for this example.
    p_dir = Path("project/specs")
    p_dir.mkdir(parents=True, exist_ok=True)
    p_file = p_dir / "materials.txt"
    p_file.touch() # .touch() creates an empty file.
    
    p_missing = Path("project/specs/missing.txt")
    
    print(f"'{p_dir.name}' exists: {p_dir.exists()}")
    print(f"'{p_dir.name}' is a directory: {p_dir.is_dir()}")
    print(f"'{p_file.name}' is a file: {p_file.is_file()}")
    print(f"'{p_missing.name}' exists: {p_missing.exists()}")
    # Outcome:
    # 'specs' exists: True
    # 'specs' is a directory: True
    # 'materials.txt' is a file: True
    # 'missing.txt' exists: False
}{Using '.exists()', '.is\_dir()', and '.is\_file()' to inspect paths.}


\hhh{Manipulating Files and Directories}
With a Path object, you can create, rename, and delete items on the file system. These actions form the core of automation scripts for project setup or cleanup.

\code{c13_tp_rename_delete}{python}{
    from pathlib import Path

    # Create a temporary file to work with.
    temp_file = Path("temp_note.txt")
    temp_file.write_text("This is a temporary file.", encoding="utf-8")

    # Define a path for an organized archive.
    archive_dir = Path("archive")
    archive_dir.mkdir(exist_ok=True) # Create the archive folder.
    archive_file = archive_dir / "final_note.txt"
    
    # Rename (or move) the file to its new location.
    # If the destination file already exists, this will raise an error.
    temp_file.rename(archive_file)
    print(f"Moved to: {archive_file}")

    # To delete a file, first check that it exists to avoid errors.
    if archive_file.exists():
        archive_file.unlink() # The .unlink() method deletes the file.
        print(f"Deleted: {archive_file}")
}{Use `.mkdir()`, `.rename()`, and `.unlink()` to manage files and folders.}

\hhh{Listing Directory Contents}
A common architectural task is to process all files within a specific project folder. The \texttt{.iterdir()} method allows you to loop through every item, while \texttt{.glob()} lets you find items matching a specific pattern.

\code{c13_tp_listing_contents}{python}{
    from pathlib import Path
    import collections

    # Setup: Create a sample project directory.
    assets_dir = Path("project_assets")
    assets_dir.mkdir(exist_ok=True)
    (assets_dir / "plan_L1.dwg").touch()
    (assets_dir / "plan_L2.dwg").touch()
    (assets_dir / "render_front.jpg").touch()
    (assets_dir / "render_side.jpg").touch()
    (assets_dir / "specs.pdf").touch()

    # Use .iterdir() to list all items and count them by type.
    file_counts = collections.defaultdict(int)
    for path in assets_dir.iterdir():
        if path.is_file():
            file_counts[path.suffix] += 1
            
    print("--- File counts by type ---")
    for suffix, count in file_counts.items():
        print(f"  {suffix}: {count} file(s)")

    # Use .glob('*.dwg') to find all drawing files specifically.
    print("\n--- Listing all drawing files ---")
    drawing_files = assets_dir.glob("*.dwg")
    for file in drawing_files:
        print(f"  Found: {file.name}")
    # Outcome:
    # --- File counts by type ---
    #   .dwg: 2 file(s)
    #   .jpg: 2 file(s)
    #   .pdf: 1 file(s)
    #
    # --- Listing all drawing files ---
    #   Found: plan_L1.dwg
    #   Found: plan_L2.dwg
}{Iterate through a directory with '.iterdir()' and find patterns with '.glob()'.}

% --------------------------------------------------------------------
\hh{Managing Project Folders with pathlib}

Projects are never just a single file. You have folders for data, drawings, reports, assets, and more. Python's \texttt{pathlib} library is the modern, intuitive, and OS-independent way to manage this file structure programmatically.

Instead of manually building strings with slashes, \texttt{pathlib} lets you create "Path" objects. The most powerful feature is using the division operator (\texttt{/}) to join path components. Python automatically uses the correct separator for the current operating system (e.g., \texttt{/} on macOS/Linux, \texttt{\\} on Windows), making your scripts robust and portable.

When creating directories with \c{.mkdir()}, two arguments are essential:
\begin{itemize}
    \item \texttt{parents=True}: Creates any missing parent directories in the path. \\
    For example, if you create \texttt{project/reports/final}, it will create \texttt{project} and \texttt{reports} if they don't already exist.
    \item \texttt{exist\_ok=True}: Prevents the script from crashing with an error if the directory you are trying to create already exists.
\end{itemize}

The following example shows a common workflow: organizing a project structure, reading data from one folder, and writing a summary report to another.

\code{c13_tp_directories}{python}{
    from pathlib import Path

    # --- 1. Define the project structure ---
    # Create Path objects for our project folders. These are not yet created on disk.
    project_root = Path("project_demo")
    data_dir = project_root / "data"      # Use the / operator to join paths
    reports_dir = project_root / "reports"  # This is much safer than "reports\\" or "reports/"

    # --- 2. Create the directories on disk ---
    # Use .mkdir() to create the actual folders.
    # The options ensure the script is safe to run multiple times.
    data_dir.mkdir(parents=True, exist_ok=True)
    reports_dir.mkdir(parents=True, exist_ok=True)

    # --- 3. Write an input file ---
    # Define the full path for a new data file inside the 'data' directory.
    input_path = data_dir / "rooms.txt"
    with open(input_path, "w", encoding="utf-8") as f:
        f.write("C-101,24.2\nC-102,20.5\nC-103,19.8\n")

    # --- 4. Process data and write a report ---
    # Define the full path for the output report inside the 'reports' directory.
    output_path = reports_dir / "summary.csv"
    total_area = 0.0
    room_count = 0

    # Open the input and output files at the same time.
    with open(input_path, "r", encoding="utf-8") as f_in, \
         open(output_path, "w", encoding="utf-8") as f_out:
        
        f_out.write("room_id,area_m2\n") # Write the header for the report
        for line in f_in:
            room_id, area_str = line.strip().split(",")
            total_area += float(area_str)
            room_count += 1
            f_out.write(f"{room_id},{area_str}\n") # Write the original data to the report

    print(f"Processed {room_count} rooms. Total area: {total_area:.2f} m2.")
    print(f"Report written to: {output_path.as_posix()}")

    # Outcome (example):
    # Processed 3 rooms. Total area: 64.50 m2.
    # Report written to: project_demo/reports/summary.csv
}{Using pathlib to create a project structure and process files.}
% ====================================================================
% PART III — PARSING STRUCTURED DATA
% ====================================================================

\hh{Parsing Common Text-Based Data Formats}

Most of the data you'll exchange between software will be in a simple text-based format. These files are lightweight, human-readable, and easy to parse. The core idea is always the same: one record per line, with values separated by a consistent character (a "delimiter").

\hhh{CSV (Comma-Separated Values)}
This is the most common format, easily exported from Excel and other software. Values in each row are separated by a comma.

\code{c13_text_csv}{python}{
    # Our data, including a header line.
    csv_data = "name,cost_per_m2\nConcrete,45.0\nSteel,82.5\nGlass,60.0\n"
    with open("materials.csv", "w", encoding="utf-8") as f:
        f.write(csv_data)

    total_cost = 0.0
    with open("materials.csv", "r", encoding="utf-8") as f:
        header = next(f) # Read and skip the first line (the header).
        for line in f:
            # .strip() removes whitespace/newlines, .split(',') breaks the line into parts.
            name, cost_str = line.strip().split(',')
            total_cost += float(cost_str)
            
    print(f"Total material cost from CSV: ${total_cost:.2f}")
    # Outcome:
    # Total material cost from CSV: $187.50
}{Parsing a CSV file by splitting each line at the comma.}

\hhh{TSV (Tab-Separated Values)}
A TSV file works exactly like a CSV, but uses a tab character (\texttt{\t}) as its delimiter. This is useful when your text data might contain commas (e.g., "Concrete, Reinforced").

\hhh{Whitespace-Delimited}
Common in scientific and engineering data (like point clouds from a survey), where columns are separated by one or more spaces. Python's \texttt{.split()} method, with no arguments, handles this automatically.

\code{c13_text_whitespace}{python}{
    # Data representing 3D coordinates.
    point_data = "x y z\n0.0 1.2 5.0\n-1.0   0.0   3.3\n" # Note the inconsistent spacing
    with open("points.txt", "w", encoding="utf-8") as f:
        f.write(point_data)

    points = []
    with open("points.txt", "r", encoding="utf-8") as f:
        next(f) # Skip header
        for line in f:
            # .split() handles any amount of whitespace between values.
            x_str, y_str, z_str = line.split()
            points.append((float(x_str), float(y_str), float(z_str)))
            
    print(f"Read {len(points)} points from whitespace-delimited file.")
    # Outcome:
    # Read 2 points from whitespace-delimited file.
}{Using .split() to parse data separated by spaces or tabs.}

% --------------------------------------------------------------------
\hh{Best Practices for Machine-Readable Data}

Before you even write a line of Python, how you structure your data files can determine if your script will be easy to write or a nightmare. A little discipline here saves hours of debugging.

\hhh{For Text Files (CSV, TSV, etc.)}
\begin{itemize}
    \item \textbf{One record per line:} Each row should represent a single, complete entry.
    \item \textbf{Use a header row:} The first line should clearly and simply name each column (e.g., \texttt{room\_id,area\_m2}). This makes your code self-documenting.
    \item \textbf{Be consistent:} Use the same delimiter and column order for every single row.
    \item \textbf{Keep numbers clean:} Store numbers as plain digits (\texttt{123.45}), not as text with units or thousand separators (\texttt{1,234.5 $m^2$}).
\end{itemize}

\hhh{For Excel Files}
Scripts work best on simple, predictable tables.
\begin{itemize}
    \item \textbf{One table per sheet:} Dedicate each sheet to a single, clean data table.
    \item \textbf{Headers in row 1:} Place column headers in the very first row.
    \item \textbf{Avoid "Human" formatting:} Do not use merged cells, blank rows to separate data, or colors to convey meaning. Your script cannot easily interpret these visual cues.
    \item \textbf{Keep data types consistent:} A column of numbers should only contain numbers, not notes like "N/A" or "pending".
\end{itemize}

% --------------------------------------------------------------------
\hh{Working with Excel Files using pandas}

Excel files (\texttt{.xlsx}) are more complex than plain text files. Manually parsing them is difficult, so we use a powerful library called \textbf{pandas}. It is the standard tool for data analysis in Python and makes reading Excel files incredibly simple. A core concept in pandas is the \textbf{DataFrame}---a smart, two-dimensional table, like a spreadsheet right inside your Python script.

\textbf{Setup:} Before running, you'll need to install the necessary libraries from your terminal:
\begin{verbatim}
pip install pandas openpyxl
\end{verbatim}

\code{c13_excel_read}{python}{
    # Pandas is conventionally imported with the alias 'pd'.
    import pandas as pd
    from pathlib import Path

    # --- 1. Create a Sample Excel File ---
    # First, create a pandas DataFrame (a data table).
    # This is like creating a spreadsheet in memory.
    df = pd.DataFrame({
        "room_id": ["E-101", "E-102", "E-103"],
        "area_m2": [30.5, 27.0, 22.5],
        "department": ["Admin", "Admin", "Studio"]
    })
    
    # Define the output path and save the DataFrame to an .xlsx file.
    # index=False prevents pandas from writing row numbers into the file.
    output_path = Path("rooms_report.xlsx")
    df.to_excel(output_path, index=False, sheet_name="RoomData")

    # --- 2. Read and Analyze the Excel File ---
    # Read the entire sheet back into a new DataFrame.
    report_df = pd.read_excel(output_path, sheet_name="RoomData")
    print("--- Full Table Preview ---")
    print(report_df.head()) # .head() shows the first 5 rows

    # Calculate the total area using built-in pandas methods.
    total_area = report_df["area_m2"].sum()
    print(f"\nTotal area from Excel: {total_area:.2f} m2")

    # It's also efficient to read only the columns you need.
    areas_only = pd.read_excel(output_path, sheet_name="RoomData", usecols=["area_m2"])
    print(f"Average area: {areas_only['area_m2'].mean():.2f} m2")
    
    # Outcome (example):
    # --- Full Table Preview ---
    #   room_id  area_m2 department
    # 0   E-101     30.5      Admin
    # 1   E-102     27.0      Admin
    # 2   E-103     22.5     Studio
    #
    # Total area from Excel: 80.00 m2
    # Average area: 26.67 m2
}{Use the pandas library to easily write and read structured data from Excel files.}

% ====================================================================
% PART IV — WRITING ROBUST AND SCALABLE SCRIPTS
% ====================================================================

Now that you have the basics of file I/O, let's look at a few common patterns. These techniques will help you write scripts that can handle large amounts of data, gracefully manage errors, and work with complex project structures. Think of these as the difference between a quick, disposable script and a powerful, reusable tool.

\hh{Processing Large Files Without Running Out of Memory}

What if your data file is gigabytes in size? Reading the entire file into memory with \texttt{.read()} or \texttt{.readlines()} will crash your script. The solution is to stream the data: read and process it one line at a time in a loop. This approach uses a tiny, constant amount of memory, regardless of the file's size.

This pattern is essential for tasks like:
\begin{itemize}
    \item Processing hourly sensor data from a building performance study.
    \item Analyzing large material schedules exported from a BIM model.
    \item Calculating statistics from massive climate data files.
\end{itemize}

\code{c13_streaming}{python}{
    from pathlib import Path

    # For demonstration, we'll create a small file. Imagine this contains
    # millions of records from a building simulation or sensor log.
    path = Path("large_sensor_log.csv")
    path.write_text("timestamp,temperature\n"
                    "2025-09-23T10:00:00,21.5\n"
                    "2025-09-23T10:01:00,21.7\n"
                    "not-a-valid-line\n"
                    "2025-09-23T10:02:00,21.8\n", encoding="utf-8")

    total_temp = 0.0
    valid_readings = 0
    with open(path, "r", encoding="utf-8") as f:
        header = next(f) # Read the header line to advance past it.
        for line in f:
            # This loop processes one line at a time, keeping memory usage low.
            parts = line.strip().split(',')
            
            # Defensive check: ensure the line has the expected number of columns.
            if len(parts) == 2:
                try:
                    # Attempt to convert the temperature to a float.
                    total_temp += float(parts[1])
                    valid_readings += 1
                except ValueError:
                    # If conversion fails (e.g., text in a number column),
                    # simply skip this malformed line and continue.
                    print(f"Skipping invalid line: {line.strip()}")
    
    avg_temp = total_temp / valid_readings if valid_readings > 0 else 0
    print(f"\nProcessed {valid_readings} valid readings.")
    print(f"Average Temperature: {avg_temp:.2f} C")
    # Outcome:
    # Skipping invalid line: not-a-valid-line
    #
    # Processed 3 valid readings.
    # Average Temperature: 21.67 C
}{Stream large files line-by-line to use constant, minimal memory.}

% --------------------------------------------------------------------
\hh{Handling Large Binary Files (Images, PDFs, Drawings)}

Binary files, like images or PDFs, must be opened in binary mode (\texttt{"rb"} or \texttt{"wb"}). Just like with large text files, you should avoid loading a huge binary file (e.g., a high-resolution render or a large drawing set) into memory all at once. The standard pattern is to copy it in smaller, fixed-size "chunks."

\code{c13_binary_chunks}{python}{
    from pathlib import Path
    
    # Imagine 'plan.pdf' is a large, multi-megabyte drawing file.
    # First, let's create a dummy source file to simulate this.
    source_path = Path("assets/plan.pdf")
    source_path.parent.mkdir(exist_ok=True)
    source_path.write_bytes(b"This is a fake PDF file." * 1000) # Dummy content

    # Define the destination for our backup copy.
    dest_path = Path("backup/plan_copy.pdf")
    dest_path.parent.mkdir(parents=True, exist_ok=True)

    # Copy the file in fixed-size chunks (e.g., 64 kilobytes).
    CHUNK_SIZE = 64 * 1024 
    
    print(f"Copying {source_path} to {dest_path}...")
    with open(source_path, "rb") as f_in, open(dest_path, "wb") as f_out:
        while True:
            # Read a chunk of data from the source file.
            block = f_in.read(CHUNK_SIZE)
            if not block:
                # An empty block means we've reached the end of the file.
                break
            # Write the chunk we just read to the destination file.
            f_out.write(block)
            
    print("Copy complete.")
}{Copy large binary files in chunks to avoid high memory consumption.}


% --------------------------------------------------------------------
\hh{Transforming Data: Reading and Writing Simultaneously}
A very common task is to read a file, modify its content in some way, and write the result to a new file. You can open both the input and output files within a single \texttt{with} block. This is cleaner and ensures both files are properly closed.

\code{c13_multiple_files_with}{python}{
    # Let's create a source file with room areas.
    with open("all_rooms.csv", "w", encoding="utf-8") as f:
        f.write("id,area_m2\n"
                "A-101,18.5\n"
                "A-102,55.1\n" # A large room
                "A-103,15.2\n"
                "A-104,72.0\n") # Another large room

    # Goal: create a new file containing only rooms larger than 50 m2.
    with open("all_rooms.csv", "r", encoding="utf-8") as fin, \
         open("large_rooms.csv", "w", encoding="utf-8") as fout:
        
        header = next(fin)
        fout.write(header) # Write the original header to the new file.
        
        for line in fin:
            room_id, area_str = line.strip().split(',')
            if float(area_str) > 50.0:
                fout.write(line) # If it meets the criteria, write the original line.

    print("Filtered report 'large_rooms.csv' created.")
    print(Path("large_rooms.csv").read_text())
    # Outcome:
    # Filtered report 'large_rooms.csv' created.
    # id,area_m2
    # A-102,55.1
    # A-104,72.0
}{Open input and output files together to create a data processing pipeline.}


% --------------------------------------------------------------------
\hh{Processing Multiple Files with {glob}}
Often, you need to perform the same action on many files at once. For example, aggregating data from several reports. \texttt{Path.glob()} lets you find all files in a directory that match a certain pattern (using wildcards like \texttt{*}).

\code{c13_globbing}{python}{
    from pathlib import Path

    # Setup: Create a directory with multiple data files from different floors.
    report_dir = Path("project/floor_reports")
    report_dir.mkdir(parents=True, exist_ok=True)
    (report_dir / "L1_areas.csv").write_text("id,area\nL1-A,100\nL1-B,150")
    (report_dir / "L2_areas.csv").write_text("id,area\nL2-A,120\nL2-B,130")
    (report_dir / "notes.txt").write_text("Project notes.")

    # Use glob("*.csv") to find all files ending with .csv in the directory.
    print(f"Finding all CSV reports in '{report_dir}'...")
    csv_files = sorted(report_dir.glob("*.csv"))
    
    grand_total_area = 0.0
    # Loop through the list of found file paths.
    for file_path in csv_files:
        print(f"  - Processing {file_path.name}...")
        with open(file_path, "r", encoding="utf-8") as f:
            next(f) # Skip header
            for line in f:
                parts = line.strip().split(',')
                if len(parts) == 2:
                    grand_total_area += float(parts[1])

    print(f"\nGrand Total Area from all reports: {grand_total_area} m2")
    # Outcome:
    # Finding all CSV reports in 'project/floor_reports'...
    #   - Processing L1_areas.csv...
    #   - Processing L2_areas.csv...
    #
    # Grand Total Area from all reports: 500.0 m2
}{Use `Path.glob()` to find and process files based on a name pattern.}


% --------------------------------------------------------------------
\hh{Defensive Programming: Handling Missing Files and Errors}

A script written for yourself might be simple, but a tool for others (or your future self) must be "defensive." It should anticipate problems like missing files or permission errors and handle them gracefully instead of crashing.

\code{c13_try_except}{python}{
    from pathlib import Path

    # Define a path to a configuration file that may or may not exist.
    config_path = Path("project/settings.conf")

    # --- First line of defense: check if the file exists. ---
    if not config_path.exists():
        print(f"Warning: Configuration file not found at '{config_path}'. Using defaults.")
    else:
        # --- Second line of defense: handle errors during file access. ---
        try:
            # This block is for code that might fail due to OS-level issues
            # (e.g., no read permissions, file is locked by another program).
            with open(config_path, "r", encoding="utf-8") as f:
                content = f.read()
                print("Configuration loaded successfully.")
        except OSError as e:
            # If any OS-related error occurs, this block will catch it.
            print(f"Error: Could not read the file '{config_path}'.")
            print(f"Reason: {e}")
            
    # Outcome (if the file is missing):
    # Warning: Configuration file not found at 'project/settings.conf'. Using defaults.
}{Write robust scripts that check for existence and catch OS-level errors.}

% ====================================================================
% MINI-LABS, PITFALLS, SUMMARY, EXERCISES
% ====================================================================

\hh{Mini-Lab — Clean a room schedule and export a report}

\textbf{Practice.} Read messy text, sanitize, compute totals, and write a clean CSV.

\code{c13_minilab_clean}{python}{
    from pathlib import Path

    src = Path("project/data/rooms_dirty.txt")
    dst = Path("project/reports/rooms_clean.csv")
    dst.parent.mkdir(parents=True, exist_ok=True)

    # Demo input to keep this lab self-contained
    with open(src, "w", encoding="utf-8") as f:
        f.write("A-401 ,  20.1\nA-402,18.9\n# comment line\n\nA-403, 21.0\n")

    count = 0
    total = 0.0

    with open(src, "r", encoding="utf-8") as f_in, open(dst, "w", encoding="utf-8") as f_out:
        f_out.write("code,area_m2\n")
        for raw in f_in:
            line = raw.strip()
            if not line or line.startswith("#"):
                continue
            parts = [p.strip() for p in line.split(",")]
            if len(parts) != 2:
                continue
            code, area_s = parts
            try:
                area = float(area_s)
            except ValueError:
                continue
            f_out.write("{},{}\n".format(code, area))
            count += 1
            total += area

    print("Rows:", count, "| Total:", round(total, 2), "| Out:", dst.as_posix())
    #Outcome:
    # Rows: 3 | Total: 60.0 | Out: project/reports/rooms_clean.csv
}{End-to-end cleanup and reporting}

% --------------------------------------------------------------------
\hh{Append-only run logs with timestamps}

\textbf{Practice.} Keep a growing log for traceability.

\code{c13_append_log}{python}{
    from pathlib import Path
    from datetime import datetime

    log_path = Path("project/reports/run.log")
    log_path.parent.mkdir(parents=True, exist_ok=True)

    stamp = datetime.now().isoformat(sep=" ", timespec="seconds")
    note = "Recomputed areas for Level 4"

    with open(log_path, "a", encoding="utf-8") as f:
        f.write("{} | {}\n".format(stamp, note))

    with open(log_path, "r", encoding="utf-8") as f:
        lines = f.read().strip().splitlines()
    print("Last entry:", lines[-1])
    #Outcome (example):
    # Last entry: 2025-09-09 12:00:00 | Recomputed areas for Level 4
}{Maintain a simple project log}

% --------------------------------------------------------------------
\hh{Beyond Manual Parsing: Python's csv and json Libraries}

So far, we've parsed simple text files by splitting each line—a great technique that works for basic data. But what happens when your data gets more complex? What if a room name contains a comma, like \texttt{"Studio, North Wing"}? Our \texttt{line.split(',')} logic would break.

To handle these real-world challenges, Python provides built-in libraries designed specifically for standard data formats. These libraries are robust, efficient, and handle all the tricky edge cases for you. The two most important are \texttt{csv} for tabular data and \texttt{json} for nested configuration data.

\hhh{csv Library: For Robust Tabular Data}
A CSV (Comma-Separated Values) file is the universal format for spreadsheet-like data. The \texttt{csv} module is Python's tool for correctly reading and writing these files. It automatically handles complexities like values that contain commas (by wrapping them in quotes) and ensures consistent row formatting. To use it, you first need to add \texttt{import csv} to the top of your script.

{Key Syntax:}
\begin{itemize}
    \item \texttt{csv.writer(file\_object)}: Creates a writer object to convert your data into delimited strings. Use its \texttt{.writerow(a\_list)} method to write a list of values as a single row.
    \item \texttt{csv.DictReader(file\_object)}: The most convenient way to read a CSV. It treats the first row as headers and returns each subsequent row as a dictionary, allowing you to access data by column name (e.g., \texttt{row["area\_m2"]}).
    \item \textbf{Note:} When opening a file for use with the \texttt{csv} module, always add the argument \texttt{newline=""} to the \texttt{open()} call. This prevents the module from writing extra blank rows on some operating systems (like Windows).
\end{itemize}

\code{c13_csv_module}{python}{
    import csv
    from pathlib import Path

    p = Path("project/data/rooms_proper.csv")
    p.parent.mkdir(parents=True, exist_ok=True)

    # Use csv.writer to automatically handle special characters like the comma in the name.
    # Note: open() includes newline="" to prevent extra blank rows.
    with open(p, "w", newline="", encoding="utf-8") as f:
        writer = csv.writer(f)
        writer.writerow(["code", "name", "area_m2"]) # Write the header row
        writer.writerow(["C-101", "Studio, North Wing", 32.8]) # This comma is handled correctly
        writer.writerow(["C-102", "Crit Space", 28.4])

    # Use csv.DictReader to read rows as dictionaries for easy data access.
    total = 0.0
    with open(p, "r", newline="", encoding="utf-8") as f:
        reader = csv.DictReader(f) # First row is automatically used as keys.
        for row in reader:
            # Access data by column name instead of by index. Much more readable!
            total += float(row["area_m2"])
            
    print(f"Total area calculated with csv.DictReader: {total:.2f} m2")
    # Outcome:
    # Total area calculated with csv.DictReader: 61.20 m2
}{Use 'csv.writer' and 'csv.DictReader' for robust tabular data handling.}

\hhh{json Library: For Settings and Nested Data}
JSON (JavaScript Object Notation) is the standard format for storing nested data, like configuration settings for a script or data from a web API. It maps directly to Python's dictionaries (for key-value pairs) and lists. The \texttt{json} library lets you effortlessly convert between these Python objects and a text-based JSON representation. To use it, start your script with \texttt{import json}.

Key Syntax:
\begin{itemize}
    \item \texttt{json.dump(python\_object, file\_object)}: Takes a Python dictionary or list and \textbf{d}umps it into a file as a JSON-formatted string. Use the argument \texttt{indent=2} to make the file neatly formatted and human-readable.
    \item \texttt{json.load(file\_object)}: Does the reverse. It \textbf{l}oads the content from a JSON file and parses it back into a Python dictionary or list.
\end{itemize}

\code{c13_json_roundtrip}{python}{
    import json
    from pathlib import Path

    # A Python dictionary is a perfect way to store nested project settings.
    project_settings = {
        "projectName": "ARC500 Design Studio",
        "units": "metric",
        "analysisParams": {"loadDL_kPa": 3.2, "loadLL_kPa": 2.4}
    }

    p = Path("project/data/settings.json")
    p.parent.mkdir(parents=True, exist_ok=True)

    # Use json.dump() to write the dictionary to a human-readable file.
    with open(p, "w", encoding="utf-8") as f:
        json.dump(project_settings, f, indent=2)

    # Use json.load() to read the file back into a Python dictionary.
    with open(p, "r", encoding="utf-8") as f:
        restored_settings = json.load(f)

    # Now you can access the data just like a normal Python dictionary.
    print(f"Project: {restored_settings['projectName']}")
    print(f"Dead Load: {restored_settings['analysisParams']['loadDL_kPa']} kPa")
    # Outcome:
    # Project: ARC500 Design Studio
    # Dead Load: 3.2 kPa
}{Use 'json.dump()' and 'json.load()' to save and load structured dictionary data.}

% --------------------------------------------------------------------
\hh{Common pitfalls and checks}

\textbf{Remember.}
\begin{itemize}
  \item \c{write()} does not add \c{"\textbackslash n"}; include it yourself.
  \item Opening with \c{"w"} erases existing content; use \c{"a"} to append.
  \item Specify \c{encoding="utf-8"} for text.
  \item For \c{csv} on Windows, pass \c{newline=""} to \c{open()}.
  \item Prefer line iteration for large files.
  \item Excel sheets: one table per sheet, header in row 1, consistent types.
\end{itemize}

% --------------------------------------------------------------------
\hh{Key Takeaways}
\begin{tcolorbox}[% passing options here
    colframe=orange,
    colback = white,
%    colback=orange!40!yellow!30, % think of mixing 2 colors with sliders
    arc=2mm % sharp corners
 ]

\begin{itemize}
  \item Safe pattern: \c{with open(path, mode, encoding="utf-8")} \dots
  \item Modes: \c{"r"}, \c{"w"}, \c{"a"}, \c{"x"} (+ \c{"b"} for binary).
  \item Reads: \c{read()}, \c{readline()}, \c{readlines()}, or iterate lines.
  \item Paths/folders: \c{pathlib.Path} (\c{exists()}, \c{mkdir()}, \c{rename()}, \c{unlink()}).
  \item Text formats: CSV/TSV/whitespace/fixed width; one record per line and a header.
  \item Excel: pandas \c{read\_excel()}, \c{usecols}, \c{sheet\_name}.
  \item Tables vs settings: \c{csv} for rows/columns; \c{json} for nested options.
\end{itemize}
\end{tcolorbox}
% --------------------------------------------------------------------
\hh{Practice exercises}

\textbf{Warm-up.} Make \texttt{materials.txt} with three lines of \texttt{name,cost\_per\_m2}. Read it, compute the average cost to two decimals, then append a new line and recompute.

\textbf{Line reader.} Create \texttt{level1\_areas.txt} with five lines of \texttt{code,area}. Write a script that prints the first line, last line, and the number of lines without loading the whole file into memory.

\textbf{Report builder.} Read \texttt{rooms\_demo.txt} and write \texttt{areas\_summary.csv} with header and totals at the bottom (count and sum).

\textbf{Excel practice.} Create \texttt{rooms.xlsx} with a sheet named \texttt{Rooms} and columns \texttt{code,area\_m2}. Read it with pandas, print row count and total area, then read only the \texttt{code} column.

\textbf{Cleanup.} Safely delete any temporary files whose names start with \texttt{tmp\_}. Hint: \texttt{Path.glob("tmp\_*.txt")} and \texttt{unlink()}.
